% !TeX program = lualatex
\documentclass[11pt]{article}

\usepackage[margin=2.5cm]{geometry}
\usepackage[T1]{fontenc}
\usepackage[scaled]{helvet}
\usepackage{pgf}
\usepackage{graphicx}
\usepackage{caption}
\usepackage{xcolor}

\newcommand{\todo}[1]{\textcolor{gray!70}{\itshape [#1]}}

\title{Electoral Pressure and Political Rhetoric:\\An Exploratory Causal Analysis with Sparse Autoencoders}
\author{Boris Ivanchikov}
\date{\today\\[0.5em]\small Draft}

\begin{document}

\maketitle

\begin{abstract}
Do politicians adapt their communication in response to changes in their electorate, and along which dimensions? We analyze the communication of U.S. House members on Twitter (now X) around the 2020 redistricting cycle, exploiting map-induced shifts in district partisanship as quasi-random variation in electoral pressure within a difference-in-differences design. Prior work quantifies rhetoric through word frequencies or pre-specified categories, constraining the scope of discoverable effects to dimensions chosen in advance. Instead, we adopt an exploratory causal inference approach, using a foundation model and a sparse autoencoder (SAE) to construct a data-driven atlas of rhetorical dimensions. Estimating treatment effects across thousands of candidate features raises a severe multiple-testing problem, which we address \todo{somehow}. We find \todo{some cool stuff}. To our knowledge, this constitutes the first application of exploratory causal inference to observational social-science data. 
\end{abstract}

\section{Introduction}

\subsection{Redistricting and Congressional Communication}

\begin{figure}[ht!]
  \centering
  {\sffamily\resizebox{\textwidth}{!}{\input{redistricting.pgf}}}
  \caption{Congressional district partisanship before and after the
  2020--2022 redistricting cycle. \todo{Full caption to be written.}}
  \label{fig:redistricting}
\end{figure}

\subsection{Sparse Autoencoders}

\todo{TODO}

\section{Approach}

\todo{TODO}

\section{Results}

\todo{TODO}

\end{document}
